% J21 — The σ²-Triadic Decomposition on Z/10Z (R1 revision 2026-05-07)
%
% Brayden R. Sanders, M. Gish, 7Site LLC
% Submitted to: Algebraic Combinatorics
% Source corpus: Atlas/LENS_TAXONOMY_2026-05-06/SIGMA2_TRIADIC_DECISION.md
%                + papers/wp_bridge_findings_2026_05_02/WP9_LATTICE_paradoxical_info_algebras.md §5
% Companion cited as already-submitted: J15
% Per-venue cap note: 2nd Algebraic Combinatorics paper after J15
%
% R1 REVISION NOTES (2026-05-07): Per fresh-eyes referee J18_AlgComb_FreshEyes.md:
%  - Proposition 4.1 (σ^2-orbit per-orbit, was Prop 5.4) statement corrected:
%      ∑_{O_1} Ψ_B = -8, ∑_{O_2} Ψ_B = -7  (proof was right; statement had typo)
%  - Ψ_B function defined inline by an explicit table in §2.3 (no longer requires
%    inaccessible companion); the values are determined uniquely by the four
%    decomposition constraints σ-cycle/σ-fixed/F/S plus T-V on the singleton roles.
%  - Linear-vs-boundary period formula tension resolved: the table is now stated
%    as the single source of truth; the period-formula heuristics are mentioned
%    only as motivation in Remark 2.1.
%  - "Conservation/manifestation duality" replaced by "two crossing decompositions
%    of a single integer-valued invariant" with the Definition 3.4 distinction
%    between table-independent and table-specific identities.
%
% MSC 2020: 05E18, 06A07, 11F06, 20N02

\documentclass[11pt,reqno]{amsart}

\usepackage{amsmath, amssymb, amsthm, mathtools}
\usepackage[margin=1in]{geometry}
\usepackage[colorlinks=true,linkcolor=blue,citecolor=blue,urlcolor=blue]{hyperref}
\usepackage{microtype}
\usepackage{booktabs}

\theoremstyle{plain}
\newtheorem{theorem}{Theorem}[section]
\newtheorem{proposition}[theorem]{Proposition}
\newtheorem{lemma}[theorem]{Lemma}
\newtheorem{corollary}[theorem]{Corollary}

\theoremstyle{definition}
\newtheorem{definition}[theorem]{Definition}
\newtheorem{conjecture}[theorem]{Conjecture}

\theoremstyle{remark}
\newtheorem*{remark}{Remark}

\newcommand{\Z}{\mathbb{Z}}
\newcommand{\F}{\mathbb{F}}
\newcommand{\TS}{\mathrm{TSML}}
\newcommand{\BH}{\mathrm{BHML}}
\newcommand{\sig}{\sigma}

\title[$\sigma^{2}$-Triadic Decomposition on $\Z/10\Z$]
{Two Crossing Decompositions of a $-21$ Invariant on $\Z/10\Z$ \\
 with the $\sigma^{2}$-Triadic Refinement}

\author{Brayden R.\ Sanders}
\address{7Site LLC, Hot Springs, Arkansas, USA}
\email{brayden@7site.co}

\author{M.\ Gish}
\address{Independent Researcher, Hot Springs, Arkansas, USA}
\email{monica.gish1992@gmail.com}

\subjclass[2020]{05E18, 06A07, 11F06, 20N02}
\keywords{Z/10Z, sigma permutation, triadic decomposition, role partition,
finite magma, triangular numbers, Fibonacci numbers}

\date{\today}

\begin{document}

\begin{abstract}
On the residue ring $\Z/10\Z$, fix the canonical order-$6$ permutation
$\sig$ with cycle structure $(0)(3)(8)(9)(1\,7\,6\,5\,4\,2)$ --- four
fixed points plus one $6$-cycle (so $\sig$ has order $\mathrm{lcm}(1,6)
= 6$; only $\sig^{3}$ is an involution). Let $\Psi_{B} : \Z/10\Z \to \Z$ be the integer-valued function
defined explicitly in Table~\ref{tab:psi-b} below.
The total $\sum_{n=0}^{9} \Psi_{B}(n) = -21$ admits two distinct
decompositions of independent combinatorial origin:
\begin{enumerate}
\item[(I)] a \emph{$\sig$-orbit decomposition} producing the triangular
split $T_5 + T_3 = 15 + 6 = 21$, with $-T_5$ summed over the $6$-cycle
and $-T_3$ summed over the four fixed points;
\item[(II)] a \emph{role-partition decomposition} producing the Fibonacci
split $F_7 + F_6 = 13 + 8 = 21$, where the role partition
$F = \{1,3,5,7,9\}$, $S = \{2,4,8\}$, $T = \{6\}$, $V = \{0\}$ comes from
the substrate operators $\TS_{8}$, $\BH_{10}$ of \cite{SandersGishFourCore},
and the singletons $T = \{6\}$ and $V = \{0\}$ contribute $-1$ and $+1$
respectively.
\end{enumerate}
The two decompositions cross: they agree on the total but partition
$\Z/10\Z$ along genuinely different congruence classes.

The $\sig^{2}$-action refines decomposition (I): $\sig^{2}$ partitions
the $6$-cycle into two triangular orbits $O_{1} = \{1, 6, 4\}$ and
$O_{2} = \{7, 5, 2\}$, with per-orbit contributions
$\sum_{O_1} \Psi_B = -8$ and $\sum_{O_2} \Psi_B = -7$, partitioning
the $\sig$-cycle component $-15$ into the indices $7$ and $8$ (the
canonical $4$-core indices $H$ and $\mathrm{Br}$) in negated form.

We give the entire ledger explicitly, in self-contained form, using
only the definitions of $\sig$, $\sig^{2}$, the role partition, and
the $\Psi_B$ table.  Every numerical claim below reduces to direct
addition of $10$ integer values, which the reader can check by
inspection.
\end{abstract}

\maketitle

\tableofcontents

%==============================================================
\section*{0. Lens, substrate, and claim tier}
\label{sec:lens}
%==============================================================

\emph{Lens and substrate.} This paper works on $\Z/10\Z$ with the
canonical order-$6$ permutation
$\sig = (0)(3)(8)(9)(1\,7\,6\,5\,4\,2)$ (only $\sig^{3}$ is an
involution) and the explicit
integer-valued function $\Psi_B : \Z/10\Z \to \Z$ tabulated in
Table~\ref{tab:psi-b}.  The table originates as a per-element
$\BH_{10}$-period contribution in the corrected $(\TS_{8},
\BH_{10})$-substrate of \cite{SandersGishFourCore,SandersBridgeWP9};
for the present paper it is treated as an input definition.  The
role partition $F \sqcup S \sqcup T \sqcup V = \Z/10\Z$ comes from the
same $\TS_{8}, \BH_{10}$-substrate.  These choices are not derived
from first principles; they reflect a structural reading of the
substrate, and the theorems below are theorems on this specific
$(\Psi_B, \sig, \text{role-partition})$ datum.  Analogous theorems
would hold on other substrate-and-table choices; whether other
substrate choices give similarly rich downstream connections is open.

\emph{Closest published neighbour.} The closest published precedent
for the family of small finite commutative non-associative structures
on $\Z/10\Z$ (and its $F_{p}$-lifts) is Drápal and
Wanless~\cite{DrapalWanless2021}, who study the opposite structural
extremum (maximally non-associative quasigroups); the present
substrate lies at the structurally regular end of the same family.

\emph{Tier discipline.}  The structural claims of this paper split as
follows.
\begin{itemize}
\item \emph{PROVED.}  Theorem~\ref{thm:sigma-orbit} (the $\sig$-orbit
triangular decomposition) is proved by direct addition from
Table~\ref{tab:psi-b}; it is independent of the choice of
substrate tables in the precise sense of Definition~\ref{def:table-spec}.
Theorem~\ref{thm:cross} (no refinement of one decomposition by the
other) is proved by direct enumeration.  Proposition~\ref{prop:per-orbit}
($\sig^{2}$-orbit per-orbit values $-8, -7$) is proved by direct
addition.
\item \emph{COMPUTED.}  Every numerical claim below reduces to a
direct sum of ten integer values from Table~\ref{tab:psi-b}.  The
companion script \texttt{verify\_J18.py} (CC-BY-4.0, standard-library
Python) checks all six load-bearing identities in under a second.
\item \emph{TABLE-DEPENDENT IDENTITIES (Definition~\ref{def:table-spec}).}
Theorem~\ref{thm:role-fib}'s $\{-F_{6}, -F_{7}\} = \{-8, -13\}$
Fibonacci split and Proposition~\ref{prop:per-orbit}'s
$\{-7, -8\}$ split on the indices $H, \mathrm{Br}$ are recorded
as empirical observations: they hold for the canonical
$\TS_{8}, \BH_{10}$ data of \cite{SandersGishFourCore} but were
broken by $0/200$ random commutative tables \cite[N8]{SandersBridgeWP9}.
\item \emph{OPEN.}  A closed-form expression for $\Psi_{B}(n)$ in $n$
alone, a forcing argument promoting the $\sig^{2}$-orbit
$\{-7, -8\}$ values from table-specific to table-independent, and the
canonical $\sig^{2}$-triadic projection of $\BH_{10}$ are recorded
in \S\ref{sec:open}.
\end{itemize}

%==============================================================
\section{Introduction}
%==============================================================

\subsection{What this paper establishes.}
We work entirely on $\Z/10\Z$ with one specified order-$6$ permutation
$\sig$, one explicitly tabulated function $\Psi_B : \Z/10\Z \to \Z$
(Table~\ref{tab:psi-b}), and one explicitly partitioned role
decomposition $F \sqcup S \sqcup T \sqcup V = \Z/10\Z$.
All numerical claims below reduce to direct sums of $10$ integer
values; the reader can check each by inspection.
The combinatorial content is in three structural facts:
\begin{enumerate}
\item the $-21$ total decomposes as $-T_5 + -T_3 = -15 + -6$ along
$\sig$-orbits (Theorem~\ref{thm:sigma-orbit}),
\item the same $-21$ total decomposes as $-F_7 + -F_6 + (-1) + 1
= -13 + -8 + -1 + 1$ along the role partition
(Theorem~\ref{thm:role-fib}), and
\item the $\sig^{2}$-refinement partitions the $\sig$-cycle component
$-15$ into the two triangular orbits with per-orbit sums
$-8$ and $-7$, the negations of the indices $\mathrm{Br}$ and $H$
($=8$ and $7$ respectively, the canonical $4$-core indices)
(Proposition~\ref{prop:per-orbit},
Theorem~\ref{thm:cm-sigma2}).
\end{enumerate}
Theorem~\ref{thm:cross} records that the two decompositions of (i) and
(ii) genuinely cross: neither refines the other, and the intersection
$F \cap \sig\text{-cycle} = \{1, 5, 7\}$ is neither a $\sig$-orbit nor
a role class.

\subsection{What this paper does \emph{not} establish.}
We are explicit about scope.
\begin{itemize}
\item We do not derive the $\Psi_B$ values from a closed-form formula
internal to this paper; they are the integer-valued data
$\{(n, \Psi_B(n)) : n = 0, \ldots, 9\}$ tabulated in
Table~\ref{tab:psi-b}, treated as a definition.  The table arose in
\cite{SandersBridgeWP9} as an integer-period contribution under
$\BH_{10}$; here it is input data.
\item We do not give a Galois-style or representation-theoretic
formalisation of the relationship between the two decompositions;
the relationship is the concrete ``two decompositions agreeing on
the total $-21$ along genuinely different congruence classes'' of
Theorem~\ref{thm:cross}, nothing more.
\item We do not commit to a canonical $\sig^{2}$-triadic projection of
$\BH_{10}$; the three candidates currently catalogued in
\cite{Atlas2026Tiering} are not promoted to canonical here.  The
$\Psi_B$-ledger we present is invariant under the choice and is the
substrate input to any future canonical promotion.
\end{itemize}

\subsection{What changed in revision.}
The pre-revision draft asserted the $\sig^{2}$-orbit per-orbit values
in Proposition~\ref{prop:per-orbit} as $\sum_{O_1} = -7$, $\sum_{O_2} =
-8$ in its statement, while its proof correctly computed $-8$ and $-7$
respectively, and the downstream ledger used the proof's values.
An external referee's independent computation confirmed the proof
(see Acknowledgments); we have corrected the proposition's statement.
The pre-revision draft also defined $\Psi_B$ by reference to a
``linear period formula'' and a ``boundary period formula'' from a
companion paper, in mutually inconsistent ways; this is now replaced
by Table~\ref{tab:psi-b} as direct input data
(\S\ref{sec:setup}, Remark~\ref{rem:periods}).
Finally, the pre-revision draft used the language of
``conservation/manifestation duality'' without a precise definition;
this is now replaced by the precise distinction
table-independent vs.\ table-specific
(Definition~\ref{def:table-spec}).

%==============================================================
\section{Setup}
\label{sec:setup}
%==============================================================

\subsection{The corrected substrate frame.}
We work on the corrected substrate frame of
\cite{SandersGishFourCore, SandersBridgeWP9}: $\TS_{8}$ acts on
$\{1, 2, 3, 4, 5, 6, 8, 9\}$ (the interior index set) and $\BH_{10}$
acts on the full $\Z/10\Z$, with the cells $V = 0$ and $H = 7$ acting
as flow boundaries.  The earlier $\TS_{10}$-frame analyses are invalid
(see \cite[N10]{SandersBridgeWP9}) and not used here.

\subsection{The permutation $\sig$ and its powers.}
We fix the permutation $\sig$ on $\Z/10\Z$ given by the cycle structure
\[
\sig \;=\; (0)(3)(8)(9)(1\,7\,6\,5\,4\,2),
\]
that is: $\sig$ has four fixed points $\{0, 3, 8, 9\}$ and one
$6$-cycle $\{1, 7, 6, 5, 4, 2\}$ traversed in the order
$1 \mapsto 7 \mapsto 6 \mapsto 5 \mapsto 4 \mapsto 2 \mapsto 1$.
The order of $\sig$ is $\mathrm{lcm}(1, 1, 1, 1, 6) = 6$; in particular
$\sig$ itself is \emph{not} an involution, $\sig^{2}$ has order $3$ on
the $6$-cycle, and only $\sig^{3}$ (the cube) is an involution on
$\Z/10\Z$.
$\sig^{2}$ has order $1$ on the four fixed points and order $3$ on the
$6$-cycle, partitioning the $6$-cycle into the two triangular orbits
\[
O_{1} \;=\; \{1, 6, 4\},
\qquad O_{2} \;=\; \{7, 5, 2\}.
\]
($O_{1}$ contains every other element of the $\sig$-walk starting at
$1$; $O_{2}$ contains every other element starting at $7$.)
$\sig$ swaps $O_1$ and $O_2$.

\subsection{The function $\Psi_B$.}
The function $\Psi_B : \Z/10\Z \to \Z$ is defined by
Table~\ref{tab:psi-b}.

\begin{table}[h]
\centering
\caption{The function $\Psi_B$.  These ten integer values are the
input data of this paper.  They originate as
$\Psi_B(n) = -(p(n) - 1) + \delta_{n, 0} + \delta_{n, 8}$ where $p(n)$
is the per-element period under $\BH_{10}$ in the substrate of
\cite{SandersGishFourCore} and the corrections at $n = 0, 8$ come
from the $\sig$-fixed period normalization
(\cite[Cor.~2.2 and Thm.~5.1]{SandersBridgeWP9}; see
Remark~\ref{rem:periods}).
For the present paper we treat Table~\ref{tab:psi-b} as input data and
do not re-prove the period values.}
\label{tab:psi-b}
\begin{tabular}{c|cccccccccc}
\toprule
$n$           & 0 & 1 & 2 & 3 & 4 & 5 & 6 & 7 & 8 & 9 \\
$\Psi_B(n)$   & $+1$ & $-5$ & $-3$ & $-2$ & $-2$ & $-1$ & $-1$ & $-3$ & $-3$ & $-2$ \\
\bottomrule
\end{tabular}
\end{table}

By direct addition,
\[
\sum_{n=0}^{9} \Psi_B(n) = 1 + (-5) + (-3) + (-2) + (-2) + (-1) + (-1) + (-3) + (-3) + (-2) = -21.
\]

\begin{remark}[On the period values]\label{rem:periods}
The pre-revision draft of this paper assigned $\Psi_B$ values via
two formulas in tension: the \emph{linear period formula}
$p(n) = 7 - n$ on the interior of the $\sig$-cycle, and the
\emph{boundary period formula} \cite[Cor.~2.2]{SandersBridgeWP9}
on the boundary cells $\{2, 6, 7\}$ of the $\sig$-cycle.
These two formulas give different values for $n \in \{2, 6, 7\}$
($7-n \neq p(n)$ at those cells), and the pre-revision proof used
the boundary values without specifying the domain restriction.
The present revision sidesteps the issue by tabulating the values
directly in Table~\ref{tab:psi-b} and treating it as the single
source of truth.  Readers interested in the $\BH_{10}$-internal
derivation can consult \cite[\S 2 and \S 5]{SandersBridgeWP9}.
\end{remark}

\subsection{The role partition.}
The functional role partition is
\[
F = \{1, 3, 5, 7, 9\},
\quad S = \{2, 4, 8\},
\quad T = \{6\},
\quad V = \{0\},
\]
with $|F| = 5$, $|S| = 3$, $|T| = 1$, $|V| = 1$.
The partition is functional in the sense of
\cite[\S 4.2]{SandersGishFourCore}: $F$ is the odd-residue class
(BHML's flow cells), $S$ is the non-zero even-residue class minus the
doubled $H$-index ($n = 7$) / $\sig$-cycle peak (BHML's structural cells),
$T = \{6\}$ is the BHML transition cell, and $V = \{0\}$ is the void.
The partition is not closed under $\TS_{8}$ or $\BH_{10}$ separately,
and crosses the $\sig$-orbit decomposition (e.g.,
$F \cap \sig\text{-cycle} = \{1, 5, 7\}$ is neither a $\sig$-orbit nor
a $\sig^{2}$-orbit).

%==============================================================
\section{The $-21$ invariant: triangular and Fibonacci splits}
\label{sec:invariant}
%==============================================================

\begin{theorem}[$\sig$-orbit triangular decomposition]\label{thm:sigma-orbit}
\[
\sum_{n \in \sig\text{-cycle}} \Psi_{B}(n)
\;=\; -T_{5}
\;=\; -15,
\qquad
\sum_{n \in \sig\text{-fixed}} \Psi_{B}(n)
\;=\; -T_{3}
\;=\; -6,
\]
totalling $-21 = -(T_5 + T_3)$.
\end{theorem}

\begin{proof}
Direct addition from Table~\ref{tab:psi-b}:
\[
\sum_{\sig\text{-cycle}} \Psi_B
\;=\;
\Psi_B(1) + \Psi_B(7) + \Psi_B(6) + \Psi_B(5) + \Psi_B(4) + \Psi_B(2)
\;=\; -5 + -3 + -1 + -1 + -2 + -3 \;=\; -15;
\]
\[
\sum_{\sig\text{-fixed}} \Psi_B
\;=\;
\Psi_B(0) + \Psi_B(3) + \Psi_B(8) + \Psi_B(9)
\;=\; +1 + -2 + -3 + -2 \;=\; -6.
\]
The triangular identities $T_5 = 15$, $T_3 = 6$ are standard.
\end{proof}

\begin{theorem}[Role-Fibonacci decomposition]\label{thm:role-fib}
\[
\sum_{n \in F} \Psi_{B}(n) \;=\; -F_{7} \;=\; -13,
\quad
\sum_{n \in S} \Psi_{B}(n) \;=\; -F_{6} \;=\; -8,
\quad
\sum_{n \in T} \Psi_B \;=\; -1,
\quad
\sum_{n \in V} \Psi_B \;=\; +1,
\]
totalling $-13 + -8 + -1 + 1 = -21$.
The Fibonacci numbers $F_7 = 13$ and $F_6 = 8$ appear on the role
classes $F$ and $S$; the singleton classes $T = \{6\}$ and $V = \{0\}$
contribute $-1$ and $+1$ respectively, cancelling at the total.
\end{theorem}

\begin{proof}
Direct addition from Table~\ref{tab:psi-b}:
\[
\sum_F \Psi_B
\;=\; \Psi_B(1) + \Psi_B(3) + \Psi_B(5) + \Psi_B(7) + \Psi_B(9)
\;=\; -5 + -2 + -1 + -3 + -2 \;=\; -13;
\]
\[
\sum_S \Psi_B
\;=\; \Psi_B(2) + \Psi_B(4) + \Psi_B(8)
\;=\; -3 + -2 + -3 \;=\; -8;
\]
\[
\sum_T \Psi_B = \Psi_B(6) = -1;
\quad
\sum_V \Psi_B = \Psi_B(0) = +1.
\]
The total $-13 - 8 - 1 + 1 = -21$ matches Theorem~\ref{thm:sigma-orbit}.
The Fibonacci identities $F_7 = 13$, $F_6 = 8$, $F_8 = 21$ are
standard.
\end{proof}

\begin{remark}[On the Fibonacci pair]
The role-Fibonacci values $(-F_7, -F_6) = (-13, -8)$ on the principal
classes $(F, S)$ are an empirical feature of the canonical
$\TS_{8}, \BH_{10}$ tables of \cite{SandersGishFourCore}: see
Definition~\ref{def:table-spec} for the precise sense in which this
identity is \emph{table-specific} (in contrast to the
table-independent $\sig$-orbit decomposition of
Theorem~\ref{thm:sigma-orbit}).
$0/200$ random commutative tables on $\Z/10\Z$ in
\cite[N8]{SandersBridgeWP9} reproduce the $(13, 8)$ split when the
role partition $F/S/T/V$ is fixed at the canonical assignment
above and a $\Psi_B$-analogue is computed from each random table's
$\BH_{10}$-analogue under the same rule.
\end{remark}

\subsection{Two crossing decompositions.}
\begin{theorem}[Two crossing decompositions]\label{thm:cross}
The two decompositions of the global $-21$ invariant agree on the total
but partition $\Z/10\Z$ along genuinely different congruence classes:
\begin{itemize}
\item $\sig$-orbit decomposition (Theorem~\ref{thm:sigma-orbit}):
$-15 + -6$ along $\sig\text{-cycle} \sqcup \sig\text{-fixed}$,
which is forced by the cycle structure of $\sig$ acting on $\Z/10\Z$.
\item Role-Fibonacci decomposition (Theorem~\ref{thm:role-fib}):
$-13 + -8 + -1 + 1$ along $F \sqcup S \sqcup T \sqcup V$,
which is table-specific to the canonical $\TS_{8}, \BH_{10}$ pair.
\end{itemize}
The intersection $F \cap \sig\text{-cycle} = \{1, 5, 7\}$ is not a
union of $\sig$-orbits nor of $\sig^{2}$-orbits: $\sig$ moves $1$ to
$7$ (both in $F$) but $\sig$ moves $7$ to $6 \in T$ (exiting $F$);
$\sig^{2}$ moves $1$ to $6 \in T$ (also exiting $F$).
Hence $\{1, 5, 7\}$ is not closed under $\sig$ or $\sig^{2}$.
Symmetrically, no $\sig$-orbit is contained in a single role class.
In particular, neither decomposition refines the other.
\end{theorem}

\begin{proof}
Theorems~\ref{thm:sigma-orbit} and \ref{thm:role-fib} give the two
decompositions.  The intersection $F \cap \sig\text{-cycle}$ is read
off the definitions: $F = \{1, 3, 5, 7, 9\}$ contains the four
odd-cycle elements $\{1, 5, 7\}$ from the $\sig$-cycle plus the
$\sig$-fixed elements $\{3, 9\}$.
The closure failures are computed directly from $\sig$ as above.
\end{proof}

\begin{definition}[Table-independent versus table-specific]
\label{def:table-spec}
Let $\Psi : \Z/10\Z \to \Z$ be a function arising from substrate
algebraic data (period, role, etc.), and let $S \subseteq \Z/10\Z$ be
a subset.  We say the identity $\sum_{n \in S} \Psi(n) = c$ is
\emph{table-independent} if its truth depends only on $\sig$-orbit
structure of $\Z/10\Z$ together with the input table of $\Psi$:
i.e., it is a sum-over-subset identity that follows from the values of
$\Psi$ alone, with $S$ a $\sig$-stable subset.
We say it is \emph{table-specific} if its truth additionally depends
on the choice of canonical $\TS_{8}, \BH_{10}$ tables: random
perturbations of those tables (with the same $\sig$ and role partition,
and with the natural $\BH_{10}$-period extension to the perturbed
table) destroy the identity at a non-trivial empirical rate.
\end{definition}

In Theorem~\ref{thm:cross}, the $\sig$-orbit decomposition is
table-independent (it follows from $\sig$-cycle structure plus
Table~\ref{tab:psi-b} addition), while the role-Fibonacci decomposition
is table-specific (the $0/200$ random-table check of
\cite[N8]{SandersBridgeWP9} is the supporting empirical evidence).
This replaces the pre-revision draft's terminology of
``conservation/manifestation duality'', which conflated the two
notions and was used as a label rather than a definition.

%==============================================================
\section{The $\sigma^{2}$-triadic refinement}
\label{sec:sigma-2-triadic}
%==============================================================

We now refine the $\sig$-orbit side of Theorem~\ref{thm:cross}
using $\sig^{2}$.

\begin{proposition}[$\sig^{2}$-orbit per-orbit contributions]\label{prop:per-orbit}
For the two $\sig^{2}$-orbits $O_{1} = \{1, 6, 4\}$ and
$O_{2} = \{7, 5, 2\}$ on the $\sig$-cycle,
\[
\sum_{n \in O_{1}} \Psi_{B}(n) \;=\; -8,
\qquad
\sum_{n \in O_{2}} \Psi_{B}(n) \;=\; -7,
\]
with $\sum_{O_1} + \sum_{O_2} = -15 = -T_5$, recovering the
$\sig$-cycle component of Theorem~\ref{thm:sigma-orbit}.
\end{proposition}

\begin{proof}
Direct addition from Table~\ref{tab:psi-b}:
\begin{align*}
O_{1} = \{1, 6, 4\}: &\quad
\Psi_{B}(1) + \Psi_{B}(6) + \Psi_{B}(4)
= -5 + -1 + -2 = -8, \\
O_{2} = \{7, 5, 2\}: &\quad
\Psi_{B}(7) + \Psi_{B}(5) + \Psi_{B}(2)
= -3 + -1 + -3 = -7.
\end{align*}
\end{proof}

\begin{remark}[On the per-orbit values]
The values $\{-8, -7\}$ negate the indices $\{8, 7\} =
\{\mathrm{Br}, H\}$, the canonical $4$-core indices in the substrate of
\cite{SandersGishFourCore}.  In the language of
Definition~\ref{def:table-spec}, this identity is table-specific:
no forcing argument internal to the $\sig$-cycle structure plus
Table~\ref{tab:psi-b} alone produces the $\mathrm{Br}/H$ values; they
arise because the $\BH_{10}$-period values that determine
Table~\ref{tab:psi-b} happen to land on $\{1, 2, 3, 5\}$ inside each
$\sig^{2}$-orbit in such a way that the per-orbit sums are integer
negatives of the canonical $4$-core indices.  See Section~\ref{sec:open}
for the natural empirical follow-up.
\end{remark}

\subsection{The $\sig^{2}$-orbit ledger.}
\begin{theorem}[$\sig^{2}$-form ledger]\label{thm:cm-sigma2}
The $\sig^{2}$-orbit decomposition of the global $-21$ invariant has
the following ledger:
\begin{center}
\begin{tabular}{l c c l}
\toprule
Component & Value & Source & Status \\
\midrule
$\sig$-fixed contribution & $-6$ & $\sum_{\{0,3,8,9\}} \Psi_B$ & table-indep. \\
$O_{1}$ contribution      & $-8$ & $\sum_{\{1,6,4\}}\Psi_B$    & table-spec.\\
$O_{2}$ contribution      & $-7$ & $\sum_{\{7,5,2\}}\Psi_B$    & table-spec.\\
\midrule
Sum                       & $-21$ & total                       & table-indep. \\
\bottomrule
\end{tabular}
\end{center}
The total $-21$ is table-independent; the $\sig^{2}$-orbit per-orbit
split into $\{-8, -7\}$ is table-specific (per
Definition~\ref{def:table-spec} and Remark above).
\end{theorem}

\begin{remark}[Lens-invariance]
Theorem~\ref{thm:cm-sigma2} uses only the $\sig^{2}$-orbit structure on
$\Z/10\Z$ plus Table~\ref{tab:psi-b}.  Both are invariant under the
choice of canonical $\sig^{2}$-triadic projection of $\BH_{10}$
\cite{Atlas2026Tiering}: the three Tier-D candidates currently
catalogued (value-rotation, index-rotation, anomaly-flip) all preserve
the ledger of Theorem~\ref{thm:cm-sigma2}.  A canonical promotion of
one such candidate would refine the ledger on the BHML side without
changing the $\Psi_B$-side data presented here.
\end{remark}

%==============================================================
\section{Status and open questions}
\label{sec:open}
%==============================================================

\subsection{Open: closed-form $\Psi_B$.}
This paper takes Table~\ref{tab:psi-b} as input data.  A closed-form
expression for $\Psi_B(n)$ in terms of $n$ alone (a single formula
covering all of $\Z/10\Z$, not the piecewise linear / boundary /
$\sig$-fixed split actually used in \cite{SandersBridgeWP9}) would
replace the table with a formula and reduce the paper to first
principles.  The piecewise nature of the $\BH_{10}$-period in
\cite{SandersBridgeWP9} suggests this may not be possible without
re-doing the substrate analysis.

\subsection{Open: $\mathrm{Br}/H$ per-orbit values.}
The per-orbit values $-8, -7$ from Proposition~\ref{prop:per-orbit}
match the canonical $4$-core indices $\mathrm{Br} = 8$ and $H = 7$ in
sign-flipped form.  Whether this is table-independent (a forced
$\sig^{2}$-orbit-internal identity given Table~\ref{tab:psi-b}
alone) or table-specific (an emergent feature of the canonical tables
in the precise sense of Definition~\ref{def:table-spec}) is open.
A natural next computation is to evaluate $\sig^{2}$-orbit sums for
random commutative tables on $\Z/10\Z$ under the rule of
\cite[\S 5]{SandersBridgeWP9}; empirical
table-specificity would distinguish the two possibilities.

\subsection{Open: role-orbit interaction.}
The role partition $F/S/T/V$ crosses the $\sig^{2}$-orbit
decomposition.  A finer ledger combining the two could reveal whether
the role-Fibonacci pattern factors through $\sig^{2}$-orbits in a
table-independent way.  This is the natural follow-up to
Theorem~\ref{thm:cm-sigma2}.

\subsection{Open: canonical $\sig^{2}$-triadic BHML.}
The framework has three Tier-D candidates for canonical
$\sig^{2}$-triadic BHML projections (value-rotation, index-rotation,
anomaly-flip) \cite{Atlas2026Tiering}.  None has a forcing argument
analogous to D78's BR-factor cancellation \cite{SandersGishFourCore}.
The decision to operate without a canonical choice is documented at
\cite{Atlas2026Tiering} Atlas SIGMA2\_TRIADIC\_DECISION; revisiting the
decision is contingent on either (i) discovering a forcing argument,
or (ii) an empirical disambiguator.

%==============================================================
\section*{Acknowledgments}
%==============================================================

This paper extracts the $\sig^{2}$-triadic content of the substrate's
$-21$ invariant in self-contained form.
We thank an anonymous referee whose careful reading uncovered
(i) the sign-swap typo in the original Proposition~\ref{prop:per-orbit}
statement (corrected here: $\sum_{O_1} = -8$, $\sum_{O_2} = -7$, in
agreement with the proof and the downstream ledger), and
(ii) the linear-vs-boundary period formula contradiction (resolved
here by tabulating $\Psi_B$ as input data in Table~\ref{tab:psi-b},
with the original two formulas mentioned only in Remark~\ref{rem:periods}).
We also drop the pre-revision draft's terminology of
``conservation/manifestation duality'' in favor of the precise
``table-independent versus table-specific'' language of
Definition~\ref{def:table-spec}; the structural content is preserved
without the ambiguous label.

All six load-bearing numerical identities of this paper are verified
at machine precision by the bundled standard-library Python script
\texttt{verify\_J18.py} (CC-BY-4.0), which checks (C1) the
$\Psi_B$-total $-21$, (C2) the $\sig$-orbit triangular split,
(C3) the role-Fibonacci split, (C4) the $\sig^{2}$-orbit per-orbit
values $-8$ and $-7$, (C5) the crossing closure failures, and
(C6) the $\sig$-permutation sanity data ($\sig^{2}$ of period 3 on
the $6$-cycle, $\sig$-swap of $O_{1}, O_{2}$).

%==============================================================
\begin{thebibliography}{99}
%==============================================================

\bibitem{SandersGishFourCore} B.~R.~Sanders, M.~Gish,
\emph{Joint Closure, Per-Coordinate Fuse Data, and a Closed-Form
Algebraic Attractor of Two Commutative Binary Operations on
$\Z/10\Z$}, submitted to \emph{Algebraic Combinatorics}, 2026 (J15).

\bibitem{SandersBridgeWP9} B.~R.~Sanders,
\emph{The LATTICE Operator and Paradoxical Information Algebras: A
Substrate-Internal Framework on $\Z/10\Z$}, submitted to
\emph{Algebra Universalis}, 2026 (J18 / WP9 Volume I).

\bibitem{Atlas2026Tiering} B.~R.~Sanders et al.,
\emph{Lens Taxonomy and Variant Catalog: Decision Records on
$\sig^{2}$-Triadic BHML Candidates},
TIG repository \texttt{Atlas/LENS\_TAXONOMY\_2026-05-06},
DOI 10.5281/zenodo.18852047, 2026.

\bibitem{TIGsynthesis2026} B.~R.~Sanders,
\emph{Trinity Infinity Geometry Synthesis},
github.com/TiredofSleep/ck (default branch),
DOI: 10.5281/zenodo.18852047, 2026.

\bibitem{DrapalWanless2021} A.~Drápal, I.~M.~Wanless,
\emph{Maximally nonassociative quasigroups},
\emph{J. Combin. Theory Ser. A} \textbf{184} (2021), 105510.

\end{thebibliography}

\end{document}
